\section{Uživatelská dokumentace}
\label{sec:dokumentace}

Tato příloha popisuje použití kompilátoru \Roth{} z příkazové řádky a základní
práci v režimu REPL.

\subsection{Požadavky}

Pro sestavení projektu je potřeba nainstalovaný Rust toolchain.
Pro spuštění některých režimů jsou dále potřeba:
\begin{itemize}
	\item \texttt{rustc} v \texttt{PATH} (používá se při \texttt{--run} a v režimu REPL),
	\item \texttt{gcc} (používá se při \texttt{--run} s backendem pro C).
\end{itemize}

\subsection{Sestavení a testy}

Projekt lze sestavit a otestovat standardními příkazy:
\begin{lstlisting}[language=bash, caption={Sestavení a testy}]
cargo build
cargo test
\end{lstlisting}

\subsection{Kompilace souboru}

Kompilace souboru probíhá spuštěním binárky \texttt{roth} se jménem vstupního
souboru. Výstupní kód je uložen do adresáře \texttt{.build/}.

\begin{lstlisting}[language=bash, caption={Kompilace souboru}]
cargo run -- path/to/program.fs
\end{lstlisting}

Volba backendu:
\begin{lstlisting}[language=bash, caption={Volba backendu}]
cargo run -- path/to/program.fs --backend rust-ir
cargo run -- path/to/program.fs --backend c-ir
\end{lstlisting}

Volitelně lze nastavit jméno výstupního souboru (stále bude umístěn do
\texttt{.build/}):
\begin{lstlisting}[language=bash, caption={Vlastní název výstupu}]
cargo run -- path/to/program.fs -o out.rs
\end{lstlisting}

\paragraph{Režim \texttt{--run}.}
Při použití \texttt{--run} kompilátor kromě vygenerování kódu automaticky spustí
překlad cílového kódu a výsledný program vykoná. Pro Rust backend se používá
\texttt{rustc -O}, pro C backend \texttt{gcc -O2}.
\begin{lstlisting}[language=bash, caption={Kompilace a spuštění}]
cargo run -- path/to/program.fs --backend rust-ir --run
cargo run -- path/to/program.fs --backend c-ir --run
\end{lstlisting}

\subsection{Ladění výpisů}

Volba \texttt{--debug} řídí množství ladicích informací:
\begin{itemize}
	\item \texttt{0} --- bez výpisů,
	\item \texttt{1} --- vypíše vygenerovaný cílový kód,
	\item \texttt{2} --- vypíše tokeny, AST, IR a statistiky optimalizací,
	\item \texttt{3} --- navíc vypíše barevně zvýrazněný cílový kód (pokud je dostupný highlighter).
\end{itemize}

Volba \texttt{--no-color} vypne barevné zvýraznění.

\subsection{REPL režim}

REPL se spustí bez argumentu vstupního souboru, nebo pomocí \texttt{-i/--interactive}:
\begin{lstlisting}[language=bash, caption={Spuštění REPL}]
cargo run
cargo run -- -i
\end{lstlisting}

V REPL lze psát běžné \Roth{} příkazy (např. \texttt{2 3 + .}) a definovat slova
pomocí \texttt{: NAME ... ;}. REPL navíc obsahuje příkazy začínající dvojtečkou
bez mezery (např. \texttt{:help}).

\begin{lstlisting}[language=bash, caption={Základní REPL příkazy}]
:help      (napoveda)
:stack     (vypis zasobniku)
:words     (seznam uzivatelskych slov)
:vars      (seznam promennych)
:clear     (vypradneni zasobniku)
:reset     (reset stavu)
:debug 2   (nastaveni debug urovne)
:quit      (ukonceni)
\end{lstlisting}

Konstrukce \texttt{INCLUDE} je podporována pouze při kompilaci ze souboru
(je implementována jako preprocesor v \texttt{src/main.rs}); v režimu REPL se
nevyhodnocuje.
